\chapter{Form-Factor Bootstrap and Local Operators}
\label{ch:integrable-form-factor-bootstrap-local-operators}

\ConventionsLorentzian

The scattering datum of Chapters~\ref{ch:integrable-scattering-analyticity-bootstrap}
and~\ref{ch:integrable-yang-baxter-internal-symmetry} is on-shell.  Local QFT
contains more data: local operator domains, Wightman distributions, and
matrix elements of local fields between finite-particle states.  In a
massive two-dimensional theory with factorized scattering, the first
operator-level datum is the family of form factors.

\section{Finite-Particle Matrix Elements}

Let \(\mathcal I\) be the set of stable particle species, \(V_a\) their
internal species spaces, and \(S_{ab}(\theta):V_a\otimes V_b\to V_b\otimes
V_a\) the two-body scattering maps.  Ordered incoming states are denoted
\[
  \ket{\theta_1,a_1;\ldots;\theta_n,a_n}_{\rm in},
  \qquad
  \theta_1>\cdots>\theta_n .
\]
For a local scalar operator \(\widehat{\mathcal O}(x)\), the \(n\)-particle
form factor is the distributional boundary value
\[
  F^{\mathcal O}_{a_1\ldots a_n}(\theta_1,\ldots,\theta_n)
  =
  \bra{\Omega}\widehat{\mathcal O}(0)
  \ket{\theta_1,a_1;\ldots;\theta_n,a_n}_{\rm in}.
\]
If the species spaces are not one-dimensional, this is a covector on
\(V_{a_1}\otimes\cdots\otimes V_{a_n}\).  Translation covariance gives
\[
  \bra{\Omega}\widehat{\mathcal O}(x)
  \ket{\theta_1,a_1;\ldots;\theta_n,a_n}_{\rm in}
  =
  \exp\!\left[-\ii x_\mu\sum_{j=1}^n p_{a_j}^\mu(\theta_j)\right]
  F^{\mathcal O}_{a_1\ldots a_n}(\theta_1,\ldots,\theta_n).
\]
Thus the nontrivial information is at \(x=0\), while the distributional
character of \(\widehat{\mathcal O}\) enters through growth and boundary-value
conditions on \(F^{\mathcal O}\).

\section{Exchange Equation}

Analytic continuation in rapidities extends the ordered form factor to
domains in \(\C^n\).  When two adjacent rapidities are exchanged, the
incoming state changes by the two-body scattering operator.  Therefore the
covector-valued form factor satisfies Watson's exchange equation
\[
  F^{\mathcal O}_{\ldots a_i a_{i+1}\ldots}
  (\ldots,\theta_i,\theta_{i+1},\ldots)
  =
  F^{\mathcal O}_{\ldots a_{i+1} a_i\ldots}
  (\ldots,\theta_{i+1},\theta_i,\ldots)
  S_{a_i a_{i+1}}(\theta_i-\theta_{i+1}),
\]
where the \(S\)-matrix acts on the two corresponding species factors before
the covector is evaluated.

\begin{proposition}[Exchange from factorized scattering]
\label{prop:form-factor-exchange}
Assume that the finite-particle scattering states satisfy asymptotic
completeness in the massive sector and that the ordered in-state convention is
related across adjacent rapidity chambers by the factorized \(S\)-matrix.
Then every local operator whose finite-particle matrix elements admit
rapidity boundary values satisfies the exchange equation above.
\end{proposition}

\begin{proof}
On a chamber \(\theta_i>\theta_{i+1}\), the displayed state is by definition
ordered.  On the neighboring chamber, the same distributional vector is
expressed in the opposite order.  Factorized scattering says that the change
of ordered basis is exactly the adjacent two-body scattering map, with all
other tensor factors unchanged.  Pairing the two expressions with
\(\bra{\Omega}\widehat{\mathcal O}(0)\) gives the stated covector identity.
\end{proof}

\section{Cyclicity, Locality, and Semi-Locality}

Moving one rapidity by \(2\pi\ii\) crosses the particle from the ket side to
the bra side and back through the local operator.  For a scalar operator
local with respect to the particle fields, the monodromy equation is
\[
  F^{\mathcal O}_{a_1a_2\ldots a_n}
  (\theta_1+2\pi\ii,\theta_2,\ldots,\theta_n)
  =
  F^{\mathcal O}_{a_2\ldots a_n a_1}
  (\theta_2,\ldots,\theta_n,\theta_1).
\]
More generally, if \(\widehat{\mathcal O}\) is semi-local with respect to a
particle of species \(a\), the right side is multiplied by a declared
monodromy phase or endomorphism \(\omega_{\mathcal O,a}\).  This semi-local
datum is part of the operator definition; it is not inferred from the
two-body \(S\)-matrix alone.

\section{Kinematic Pole Equation}

Let \(\bar a\) denote the antiparticle species of \(a\).  The pair
\((\bar a,a)\) at rapidity difference \(\ii\pi\) has total momentum zero and
can annihilate into the vacuum.  Consequently a local form factor may have a
kinematic pole at
\[
  \theta_{\bar a}=\theta_a+\ii\pi .
\]
For a scalar local operator the residue has the structural form
\[
  -\ii\,\operatorname*{Res}_{\theta'=\theta+\ii\pi}
  F^{\mathcal O}_{\bar a\,a\,a_1\ldots a_n}
  (\theta',\theta,\theta_1,\ldots,\theta_n)
  =
  \left(
    1-\prod_{j=1}^n
    S_{a a_j}(\theta-\theta_j)
  \right)
  F^{\mathcal O}_{a_1\ldots a_n}(\theta_1,\ldots,\theta_n),
\]
with tensor contractions and charge-conjugation maps inserted according to
the species spaces.  The product is ordered by moving the particle \(a\)
through the \(n\) remaining particles before annihilation.

\begin{proof}[Derivation of the residue structure]
Insert a complete finite-particle resolution of the identity between
\(\widehat{\mathcal O}(0)\) and the adjacent pair.  Near
\(\theta'-\theta=\ii\pi\), the two-particle phase-space denominator develops
the one-particle-vacuum annihilation singularity.  If the pair annihilates
before crossing the spectator particles, the contribution is the first term.
If the particle is first moved through all spectators using factorized
scattering and then annihilates, the contribution is the second term.  Local
commutativity fixes the relative sign.  The residue is therefore the
difference of these two orderings.  A fully rigorous theorem requires a
specified Hilbert-space construction of the meromorphic finite-particle
matrix elements; the displayed equation is the form-factor axiom used after
that construction has supplied the required domains and boundary values.
\end{proof}

\section{Free Majorana Examples}

The form-factor axioms become concrete already in the massive Ising field
theory.  In the fermionic particle basis there is one self-conjugate particle
of mass \(m\), and the diagonal two-body scattering function is the constant
\[
  S(\theta)=-1 .
\]
The constant phase is physically important: the ordered rapidity variables
are still bosonic coordinates on asymptotic particle states, but exchanging
two particles multiplies the finite-particle basis by \(-1\).  Thus a scalar
form factor must be antisymmetric in rapidity variables when it has an even
number of Majorana particles.

\begin{example}[Energy-density form factor]
\label{ex:free-majorana-energy-form-factor}
Fix a scalar normalization constant \(\kappa_{\varepsilon}\).  The massive
Majorana energy-density operator \(\varepsilon\) may be normalized by
\[
  F_0^\varepsilon=0,\qquad
  F_2^\varepsilon(\theta_1,\theta_2)
  =
  \kappa_{\varepsilon}
  \sinh\!\left(\frac{\theta_1-\theta_2}{2}\right),
  \qquad
  F_n^\varepsilon=0\quad(n\neq2).
\]
The exchange equation follows because
\[
  \sinh\!\left(\frac{\theta_2-\theta_1}{2}\right)(-1)
  =
  \sinh\!\left(\frac{\theta_1-\theta_2}{2}\right).
\]
Cyclicity follows from
\[
  \sinh\!\left(\frac{\theta_1+2\pi\ii-\theta_2}{2}\right)
  =
  -\sinh\!\left(\frac{\theta_1-\theta_2}{2}\right)
  =
  \sinh\!\left(\frac{\theta_2-\theta_1}{2}\right).
\]
The kinematic-pole equation is also satisfied: the displayed form factors
are entire, and the right side vanishes because \(F_0^\varepsilon=0\) for no
spectators and because \(1-S^2=0\) for two spectators.  This example is
finite rather than recursive; it checks that the sign conventions of Watson
exchange and \(2\pi\ii\) cyclicity are compatible with a familiar local
quadratic field.
\end{example}

The order/disorder fields of the Ising theory are more instructive because
their form factors form an infinite family.  We record the odd family in a
normalization \(v\in\C\):
\begin{equation}
  F_{2k+1}^{\Sigma}(\theta_1,\ldots,\theta_{2k+1})
  =
  v\,\ii^k
  \prod_{1\le i<j\le 2k+1}
  \tanh\!\left(\frac{\theta_i-\theta_j}{2}\right),
  \qquad
  F_{2k}^{\Sigma}=0 .
  \label{eq:ising-order-form-factors}
\end{equation}
The constant \(v=F_1^\Sigma\) is the one-particle normalization in this
sector.  The complementary even family is obtained by exchanging the order
and disorder sectors; which of the two is called \(\sigma\) or \(\mu\)
depends on the sign convention for the massive Ising perturbation.  The
mathematical point here is independent of that naming convention.

\begin{proposition}[Ising odd form-factor family]
\label{prop:ising-odd-form-factor-family}
The family \eqref{eq:ising-order-form-factors} satisfies Watson exchange,
the local cyclicity equation for odd particle number, and the kinematic-pole
equation with \(S=-1\):
\[
  -\ii\,\operatorname*{Res}_{\theta'=\theta+\ii\pi}
  F_{n+2}^{\Sigma}(\theta',\theta,\eta_1,\ldots,\eta_n)
  =
  \left(1-(-1)^n\right)F_n^\Sigma(\eta_1,\ldots,\eta_n)
\]
for every odd \(n\).  For even \(n\), both sides vanish because
\(F_n^\Sigma=F_{n+2}^\Sigma=0\).
\end{proposition}

\begin{proof}
Set
\[
  P_N(\theta_1,\ldots,\theta_N)
  =
  \prod_{1\le i<j\le N}
  \tanh\!\left(\frac{\theta_i-\theta_j}{2}\right).
\]
Interchanging adjacent variables changes exactly one factor by a sign and
permutes the rest, so \(P_N\) is antisymmetric.  Since \(S=-1\), this is
precisely Watson's equation.

For \(N\) odd,
\[
  P_N(\theta_1+2\pi\ii,\theta_2,\ldots,\theta_N)
  =
  P_N(\theta_1,\theta_2,\ldots,\theta_N),
\]
because \(\tanh(z+\pi\ii)=\tanh z\).  Moving \(\theta_1\) from the first
slot to the last reverses its order relative to \(N-1\) variables and hence
multiplies the antisymmetric product by \((-1)^{N-1}=1\).  This proves
cyclicity for the odd family.

It remains to compute the annihilation residue.  Put
\[
  \theta'=\theta+\ii\pi+\epsilon .
\]
The pair factor is
\[
  \tanh\!\left(\frac{\theta'-\theta}{2}\right)
  =
  \tanh\!\left(\frac{\ii\pi+\epsilon}{2}\right)
  =
  \coth\!\left(\frac{\epsilon}{2}\right)
  =
  \frac{2}{\epsilon}+O(\epsilon).
\]
For each spectator rapidity \(\eta_j\),
\[
  \tanh\!\left(\frac{\theta'-\eta_j}{2}\right)
  \tanh\!\left(\frac{\theta-\eta_j}{2}\right)
  \longrightarrow
  \coth\!\left(\frac{\theta-\eta_j}{2}\right)
  \tanh\!\left(\frac{\theta-\eta_j}{2}\right)
  =
  1 .
\]
Therefore
\[
  \operatorname*{Res}_{\theta'=\theta+\ii\pi}
  P_{n+2}(\theta',\theta,\eta_1,\ldots,\eta_n)
  =
  2P_n(\eta_1,\ldots,\eta_n).
\]
If \(n=2k+1\), then \(F_n^\Sigma=v\ii^kP_n\), while
\(F_{n+2}^\Sigma=v\ii^{k+1}P_{n+2}\).  Hence
\[
  -\ii\,\operatorname*{Res} F_{n+2}^\Sigma
  =
  -\ii\,2v\ii^{k+1}P_n
  =
  2v\ii^kP_n
  =
  2F_n^\Sigma .
\]
Since \(n\) is odd, \(1-(-1)^n=2\), giving the claimed kinematic-pole
equation.
\end{proof}

\section{Bound-State Poles and Operator Couplings}

Suppose species \(a\) and \(b\) have a bound state \(c\).  If the
two-body \(S\)-matrix has a physical-strip pole at \(\theta=\ii u_{ab}^c\)
with residue tensor \(\Gamma_{ab}^c:V_c\to V_a\otimes V_b\), then the form
factor has the compatible pole
\[
  -\ii\,\operatorname*{Res}_{\theta_a-\theta_b=\ii u_{ab}^c}
  F^{\mathcal O}_{a b a_1\ldots a_n}
  =
  F^{\mathcal O}_{c a_1\ldots a_n}\,\Gamma_{ab}^c ,
\]
again with the stated tensor-ordering convention.  This equation is the
operator-level statement that a local field couples to the bound state through
the same one-particle channel that appears in the scattering residue.

\section{Form-Factor Reconstruction Boundary}

\begin{definition}[Form-factor datum for a local operator]
\label{def:form-factor-datum-local-operator}
A form-factor datum for a scalar local operator consists of meromorphic
covectors \(F_n^{\mathcal O}\) for all \(n\geq0\), species labels, and
semi-locality monodromy data, satisfying exchange, cyclicity, kinematic-pole,
bound-state-pole, Lorentz covariance, and specified growth conditions.
\end{definition}

The growth condition is essential.  It controls whether the spectral series
for Wightman functions,
\[
  \langle\Omega|\widehat{\mathcal O}(x)\widehat{\mathcal O}(0)|\Omega\rangle
  =
  \sum_{n=0}^{\infty}\frac{1}{n!}
  \sum_{a_1,\ldots,a_n}
  \int d\mu_n\,
  e^{-\ii x\cdot\sum_jp_j}
  F^{\mathcal O}_{a_1\ldots a_n}(\theta)
  \overline{F^{\mathcal O}_{a_1\ldots a_n}(\theta)}
\]
defines a distribution.  Here \(d\mu_n\) is the Lorentz-invariant
\(n\)-particle rapidity measure with the normalization fixed by the
one-particle Hilbert-space convention.

\begin{openproblem}[Local reconstruction from form factors]
\label{op:form-factor-local-reconstruction}
Given a factorized \(S\)-matrix and a family of form-factor data satisfying
the displayed equations and explicit growth bounds, construct the Wightman
domain, prove convergence of the spectral series as distributions, verify
local commutativity, and identify a dense set of local operators.  This is
the operator-level completion of the scattering bootstrap.
\end{openproblem}
